\section{Detectability for a physical projection chain}

\begin{definition}[The physical even and odd products]
    \label{def:physical_detectability}
    \lean{TNLean.MPS.AreaLaw.physicalParityProduct}
    \lean{TNLean.MPS.AreaLaw.physicalDetectability}
    \leanok
    Let $h_i$ be orthogonal projections supported on the nearest-neighbour
    bonds of a finite open chain, numbered from one, and put $H=\sum_i h_i$.
    Define
    \begin{align*}
        P_{\mathrm{odd}}&=\prod_{i\text{ odd}}(\Id-h_i),&
        P_{\mathrm{even}}&=\prod_{i\text{ even}}(\Id-h_i),&
        \operatorname{DL}(H)&=P_{\mathrm{even}}P_{\mathrm{odd}}.
    \end{align*}
    The terms in each product commute, and empty products are the identity.
    This is the operator in \cite[the definition of the one-dimensional
    detectability operator]{Arad2016RigorousRG}, after the source's reduction
    to projection terms.
\end{definition}

\begin{lemma}[The parity products are projections]
    \label{thm:physical_parity_projection}
    \lean{TNLean.MPS.AreaLaw.physicalParityProduct_isStarProjection}
    \leanok
    Both $P_{\mathrm{odd}}$ and $P_{\mathrm{even}}$ are orthogonal projections.
\end{lemma}
\begin{proof}\leanok
    \uses{def:physical_detectability}
    Distinct bonds of the same parity are disjoint. Operators supported on
    disjoint sites commute, so each product is a product of commuting
    orthogonal projections.
\end{proof}

\begin{lemma}[The physical detectability energy estimate]
    \label{thm:physical_detectability_energy}
    \lean{TNLean.MPS.AreaLaw.re_inner_physicalDetectability_le}
    \leanok
    For every vector $\psi$, put $\phi=\operatorname{DL}(H)\psi$. Then
    \begin{align*}
        \langle\phi,H\phi\rangle
        \le 4\bigl(\|\psi\|^2-\|\phi\|^2\bigr).
    \end{align*}
\end{lemma}
\begin{proof}\leanok
    \uses{def:physical_detectability}
    Enumerate the actual odd terms followed by the actual even terms.
    A term can fail to commute only with its two adjacent bonds.
    The detectability lemma of \cite{Arad2016RigorousRG}, in the general
    Hilbert-space form proved in \cite{QICLean2026Blueprint}, therefore
    applies with noncommutation count two. The sum and product in this
    ordering are exactly $H$ and $\operatorname{DL}(H)$.
\end{proof}

\begin{lemma}[The physical converse estimate]
    \label{thm:physical_detectability_converse}
    \lean{TNLean.MPS.AreaLaw.one_sub_three_mul_le_norm_sq_physicalDetectability}
    \leanok
    Every unit vector $\psi$ satisfies
    $\|\operatorname{DL}(H)\psi\|^2\ge1-3\langle\psi,H\psi\rangle$.
\end{lemma}
\begin{proof}\leanok
    \uses{thm:physical_parity_projection}
    Apply the converse estimate for two commuting projection families
    from \cite{QICLean2026Blueprint} to the actual even and odd terms.
    Their sums add to $H$. This strengthens the constant four and the
    eigenvector restriction in the stated converse of
    \cite{Arad2016RigorousRG}.
\end{proof}

\begin{lemma}[The full ground kernel and its orthogonal complement]
    \label{thm:physical_detectability_kernel}
    \lean{TNLean.MPS.AreaLaw.physicalParityProduct_fixes_kernel}
    \lean{TNLean.MPS.AreaLaw.physicalDetectability_fixes_kernel}
    \lean{TNLean.MPS.AreaLaw.inner_physicalDetectability_eq}
    \leanok
    Both parity products and $\operatorname{DL}(H)$ fix every vector
    $z\in\ker H$. For every such $z$ and every $\psi$,
    $\langle z,\operatorname{DL}(H)\psi\rangle=\langle z,\psi\rangle$.
    In particular $\operatorname{DL}(H)$ preserves $(\ker H)^\perp$.
    No assumption on the dimension of $\ker H$ is required.
\end{lemma}
\begin{proof}\leanok
    \uses{thm:physical_parity_projection}
    The identity $\langle z,Hz\rangle=\sum_i\|h_i z\|^2$ implies
    $h_i z=0$ for each bond. Thus every complementary factor fixes $z$.
    Self-adjointness of each parity product gives the inner-product identity.
\end{proof}

\begin{lemma}[Contraction above a unique frustration-free ground state]
    \label{thm:physical_detectability_gap}
    \lean{TNLean.MPS.AreaLaw.norm_sq_physicalDetectability_le_of_gappedGroundState}
    \leanok
    Let $\Gamma$ be a unit vector with $H\Gamma=0$, and suppose
    $H\ge\gamma(\Id-|\Gamma\rangle\langle\Gamma|)$ for $\gamma>0$.
    Every unit vector $\psi\perp\Gamma$ satisfies
    $\|\operatorname{DL}(H)\psi\|^2\le4/(4+\gamma)$.
\end{lemma}
\begin{proof}\leanok
    \uses{thm:physical_detectability_energy,thm:physical_detectability_kernel}
    Set $\phi=\operatorname{DL}(H)\psi$. The kernel identity gives
    $\phi\perp\Gamma$. Combining the actual gap inequality with the
    energy estimate yields
    $\gamma\|\phi\|^2\le\langle\phi,H\phi\rangle
      \le4(1-\|\phi\|^2)$.
    Rearranging gives the claimed squared-norm bound, without any upper
    bound on $\gamma$.
\end{proof}

\begin{lemma}[Contraction above the full zero-energy kernel]
    \label{thm:physical_detectability_spectral_gap}
    \lean{TNLean.MPS.AreaLaw.norm_sq_physicalDetectability_le_of_spectral_gap}
    \leanok
    Suppose $\gamma>0$ and
    $\operatorname{spec}(H)\subseteq\{0\}\cup[\gamma,\infty)$.
    Every $\psi\perp\ker H$ satisfies
    \begin{align}
        \|\operatorname{DL}(H)\psi\|^2
          \leq\frac4{4+\gamma}\|\psi\|^2.
        \label{eq:physical_detectability_spectral_gap}
    \end{align}
    The kernel may have arbitrary dimension, and no upper bound on
    $\gamma$ is imposed.
\end{lemma}
\begin{proof}\leanok
    \uses{thm:physical_detectability_energy,thm:physical_detectability_kernel}
    Put $\phi=\operatorname{DL}(H)\psi$. Kernel invariance gives
    $\phi\perp\ker H$. The spectral decomposition and the energy estimate give
    \begin{align}
        \gamma\|\phi\|^2\leq\langle\phi,H\phi\rangle
          \leq4(\|\psi\|^2-\|\phi\|^2).
    \end{align}
    Rearranging proves~(\ref{eq:physical_detectability_spectral_gap}).
\end{proof}

\begin{lemma}[Support of a range projection]
    \label{thm:physical_range_projection_support}
    \lean{QuantumCircuit.supportProj_mem_supportedOperators}
    \leanok
    Let the local dimension be positive. If a Hermitian operator $A$ acts on
    a set $K$ of sites, its range projection $P_A$ acts on $K$ as well.
\end{lemma}
\begin{proof}\leanok
    Every operator acting outside $K$ commutes with $A$ and hence with every
    real function of $A$. The commutant characterization of support applies
    to $P_A=\mathbf1_{\mathbb R\setminus\{0\}}(A)$.
\end{proof}

\begin{lemma}[Detectability for positive contraction chains]
    \label{thm:physical_range_detectability}
    \lean{TNLean.MPS.AreaLaw.norm_sq_physicalDetectability_supportProj_le}
    \leanok
    Let $H=\sum_i h_i$ be a nearest-neighbour chain of positive contractions
    in positive local dimension. Suppose every nonzero eigenvalue of $H$ is
    at least $\gamma>0$. Set $P_i=P_{h_i}$. Then every
    $x\perp\ker H$ satisfies
    \[
      \|\operatorname{DL}(P)x\|^2
          \le\frac4{4+\gamma}\|x\|^2.
    \]
    No condition on the dimension of $\ker H$ is required.
\end{lemma}
\begin{proof}\leanok
    \uses{thm:physical_range_projection_support,thm:physical_detectability_kernel,
      thm:physical_detectability_energy}
    The range projections have the original bond supports. Their sum dominates
    $H$ and has exactly the same kernel. The even/odd product preserves the
    orthogonal complement of this kernel. For $y=\operatorname{DL}(P)x$,
    \[
      \gamma\|y\|^2\le\langle y,Hy\rangle
        \le\Big\langle y,\sum_iP_i y\Big\rangle
        \le4\bigl(\|x\|^2-\|y\|^2\bigr).
    \]
    Rearranging gives the claim. This implements the range-projection step
    preceding the detectability argument of
    \cite[Section 6.1]{Arad2016RigorousRG}.
\end{proof}

The regional truncation in \cite{Arad2016RigorousRG} also requires a rank
correction: its printed bound $d^2+2$ fails on an eight-qubit chain with
unique ground state and gap one. The corrected bound is $d^2+3$, as
shown in \cite{gap:alvv16_frustration_free_truncation_rank}. The regional
rank theorem is separate from the detectability results above.

\begin{lemma}[Fixed vectors of the physical parity product]
    \label{thm:physical_parity_product_fixed_space}
    \lean{TNLean.MPS.AreaLaw.physicalParityProduct_apply_eq_self_iff}
    \leanok
    Let the bond operators $h_i$ be orthogonal projections. Fix $p\in\mathbb N$
    and select the indices for which $i\bmod2=p$. The actual parity product
    fixes $x$ if and only if $h_ix=0$ for every selected bond.
    This equivalence needs neither commutativity nor locality.
\end{lemma}
\begin{proof}\leanok
    \uses{def:physical_detectability}
    Extend the selected bond sequence by zero terms and apply the
    common-kernel characterization of complementary projection products
    from \cite{QICLean2026Blueprint}. All padding terms vanish.
\end{proof}

\begin{definition}[The regional frustration-free truncation]
    \label{def:physical_ff_regional_truncation}
    \lean{TNLean.MPS.AreaLaw.intervalBondRestriction}
    \lean{TNLean.MPS.AreaLaw.physicalFrustrationFreeTruncation}
    \leanok
    Let $h_i$ be orthogonal projections supported on consecutive bonds
    of a finite chain. For a site interval $J=[a,b)$, put
    $h_i^J=h_i$ when $a\le i$ and $i+1<b$, and $h_i^J=0$ otherwise.
    Let $Q_p$ be the product of $I-h_i^J$ over the indices
    $i\bmod2=p$, in the order of Definition~\ref{def:physical_detectability}.
    Define
    \[
        \widetilde H_J=
        \sum_{i:\,\{i,i+1\}\not\subseteq J}h_i+(I-Q_0)+(I-Q_1).
    \]
    This is the regional construction of
    \cite[Section 6.1]{Arad2016RigorousRG}.
\end{definition}

\begin{lemma}[Positivity and the full kernel of the regional truncation]
    \label{thm:physical_ff_regional_kernel}
    \lean{TNLean.MPS.AreaLaw.physicalFrustrationFreeTruncation_posSemidef}
    \lean{TNLean.MPS.AreaLaw.physicalFrustrationFreeTruncation_kernel_iff}
    \lean{TNLean.MPS.AreaLaw.norm_physicalFrustrationFreeTruncation_region_le_two}
    \leanok
    With the preceding definitions,
    \[
        \widetilde H_J\ge0,\qquad
        \ker\widetilde H_J=\ker\sum_i h_i,\qquad
        \|(I-Q_0)+(I-Q_1)\|\le2.
    \]
    The interval may be empty or touch a chain endpoint.
    No ground-state uniqueness or spectral gap is required.
\end{lemma}
\begin{proof}\leanok
    \uses{def:physical_ff_regional_truncation,thm:physical_parity_product_fixed_space}
    The bonds in each parity are disjoint, so $Q_0$ and $Q_1$ are
    orthogonal projections. Hence all displayed summands are positive,
    and each regional complement has norm at most one.
    A sum of orthogonal projections annihilates a vector precisely when
    every summand does. For the two regional complements, this is
    equivalent to $Q_0x=Q_1x=x$, which in turn is equivalent to
    $h_i^Jx=0$ for every internal bond. Together with the outside bonds,
    these are exactly the common-kernel conditions for the original sum.
\end{proof}

\begin{lemma}[Partition and removed mass of the physical parity products]
    \label{lem:physical_parity_partition_mass}
    \lean{TNLean.MPS.AreaLaw.physicalParityProduct_partition}
    \lean{TNLean.MPS.AreaLaw.sum_norm_sq_sub_physicalParityProduct_le}
    \leanok
    Let $P_p$ denote the full complementary projection product of parity
    $p\in\{0,1\}$, and let $H=\sum_i h_i$. Any partition of the bonds
    splits $P_p$ into its two restricted parity products. Moreover, for
    every vector $x$,
    \[
        \|(I-P_0)x\|^2+\|(I-P_1)x\|^2\le\langle x,Hx\rangle.
    \]
\end{lemma}
\begin{proof}\leanok
    \uses{def:physical_detectability}
    Bonds of the same parity have disjoint support and commute. Thus
    their complementary projections may be regrouped. The commuting
    projection union bound controls each removed mass by the sum of
    the corresponding bond energies. Adding the two parities gives
    the stated inequality.
\end{proof}

\begin{lemma}[The original detectability product and truncated energy]
    \label{lem:physical_ff_truncated_converse}
    \lean{TNLean.MPS.AreaLaw.sum_norm_sq_sub_physicalParityProduct_le_truncation}
    \lean{TNLean.MPS.AreaLaw.one_sub_three_mul_truncated_energy_le_norm_sq_physicalDetectability}
    \leanok
    Let $\widetilde H_J$ be the actual regional truncation above.
    For every vector $x$,
    \[
        \|(I-P_0)x\|^2+\|(I-P_1)x\|^2
        \le\langle x,\widetilde H_Jx\rangle.
    \]
    Consequently, for every unit vector $x$,
    \[
        1-3\langle x,\widetilde H_Jx\rangle
        \le\|P_1P_0x\|^2.
    \]
\end{lemma}
\begin{proof}\leanok
    \uses{def:physical_ff_regional_truncation,lem:physical_parity_partition_mass}
    Write each full parity product as the internal product followed by
    the outside product. For an orthogonal projection $Q$,
    $\|x-Qy\|^2\le\|x-Qx\|^2+\|x-y\|^2$.
    Apply this inequality to the internal parity products and bound
    the outside removed masses by the outside bond energies.
    The converse estimate for two orthogonal projections then gives
    the second inequality.
\end{proof}

\begin{lemma}[The nonzero gap of the regional frustration-free truncation]
    \label{thm:physical_ff_regional_gap}
    \lean{TNLean.MPS.AreaLaw.physicalFrustrationFreeTruncation_energy_ge_of_spectral_gap}
    \lean{TNLean.MPS.AreaLaw.physicalFrustrationFreeTruncation_spectral_gap}
    \leanok
    Suppose that $\gamma>0$ and
    $\operatorname{spec}(H)\subseteq\{0\}\cup[\gamma,\infty)$.
    Then
    \[
        \operatorname{spec}(\widetilde H_J)
        \subseteq\{0\}\cup
        \left[\frac{\gamma}{3(4+\gamma)},\infty\right).
    \]
    Every unit vector orthogonal to $\ker\widetilde H_J$ has energy
    at least $\gamma/[3(4+\gamma)]$.
    No uniqueness assumption or upper bound on $\gamma$ is needed.
    In particular, for $0<\gamma\le1$ the lower bound is at least
    $\gamma/15$. This proves the gap conclusion of
    \cite[Section 6.1]{Arad2016RigorousRG}; the false printed
    additive rank constant remains separate.
\end{lemma}
\begin{proof}\leanok
    \uses{thm:physical_ff_regional_kernel,lem:physical_ff_truncated_converse}
    The kernels of $H$ and $\widetilde H_J$ agree. On their common
    orthogonal complement, detectability gives
    $\|P_1P_0x\|^2\le4/(4+\gamma)$ for unit $x$.
    Combining this with the preceding converse estimate and rearranging
    yields the energy bound. Apply that bound to every nonzero
    eigenvector of the positive operator $\widetilde H_J$ to obtain
    the spectral conclusion.
\end{proof}
